
\begin{theorem}  \label{col:fubini} \textbf{Teorema de Fubini} en $\mathbb{R}^{2}.$  Sea $f:D\to\R$ continua sobre  $D\subseteq\Rn{2}$  un conjunto elemental   (\ref{conj_elem}) 
   
\begin{enumerate} 
        \item  Si $D$ es de Tipo 1 (\ref{def:tipo1})   entonces la integral doble de $f$ sobre $D$ se puede calcular
        \[
            \iint_D f=\int_a^b\left(\int_{\phi_1(x)}^{\phi_2(x)}f(x,y)\:dy\right)dx.  
        \]
       
        \item   Si $D$ es de Tipo 2 (\ref{def:tipo2})  entonces la integral doble de $f$ sobre $D$ se puede calcular
        \[
            \iint_D f=\int_c^d\left(\int_{\psi_1(y)}^{\psi_2(y)}f(x,y)\:dx\right)dy.  
        \]
       
        \item Si $D$ es de Tipo 3   entonces la integral doble de $f$ sobre $D$ se puede calcular  
         \[
            \iint_D f=\int_a^b\left(\int_{\phi_1(x)}^{\phi_2(x)}f(x,y)\:dy\right)dx=\int_c^d\left(\int_{\psi_1(y)}^{\psi_2(y)}f(x,y)\:dx\right)dy. 
        \]
    \end{enumerate}
\end{theorem}


\begin{obs} La \'unica manera de demostrar que un conjunto es elemental es dando su descripci\'on  impl\'icita.
\end{obs}


\begin{example}
     Sean $f(x,y)=x^2y$ y  $D$ es la regi\'on del plano encerrada por $y=x^3$, $y=x^2$  con $x\in[0,1]$,  calcular,  $$\iint_D f \:dA.$$  
\end{example}

\begin{figure}[htbp]
    \centering
    % Llamamos al archivo con el código TikZ extraído
    \begin{center}
    \begin{tikzpicture}
        \begin{axis}[
            axis lines = middle,
            xmin = 0, xmax = 1.5,
            ymin = 0, ymax = 1.5,
            xlabel = {$x$},
            ylabel = {$y$},
            domain = 0:1.25,
            samples = 100,
            xtick={0.5,1},
            ytick={0.5,1},
            clip=false,
            ]
        
            % Define the curves
            \addplot[name path=A, blue, thick, domain=0:1.22] {x^2} node[pos=0.75, right] {$y=x^2$};
            \addplot[name path=B, red, thick, domain=0:1.15] {x^3} node[pos=0.3, below right] {$y=x^3$};
        
            % Fill the area between the curves
            \addplot[orange!50, fill opacity=0.5] fill between[of=A and B, soft clip={domain=0:1}];
        
            % Label the enclosed area as D
            \node at (axis cs:0.6,0.29) {$D$};
        
        \end{axis}
    \end{tikzpicture}
    \end{center}
 
    \label{fig:ejemplo 812}
\end{figure}

\textbf{Soluci\'on:} $D$ puede ser descripto de la siguiente manera
    \[
        D=\{(x,y)\in\Rn{2}:0\leq x\leq1;\;x^3\leq y\leq x^2\},  
    \]
 por lo tanto $D$ es elemetal de Tipo 1 y  por el \autoref{col:fubini}
    \begin{gather*}
        \iint_D x^2y\:dA=\int_0^1\left(\int_{x^3}^{x^2}\:dy\right)dx=\int_0^1\left(x^2\frac{1}{2}y^2\Big\lvert_{x^3}^{x^2}\right)dx\\
        =\int_0^1\frac{1}{2}x^2\left((x^2)^2-(x^3)^2\right)=\int_0^1\frac{1}{2}x^6-\frac{1}{2}x^8\;dx\\=\frac{1}{14}x^7-\frac{1}{18}x^9\Big\lvert_0^1=\frac{1}{14}-\frac{1}{18}=\frac{1}{63}.
    \end{gather*}


\begin{tarea}Calcular  $$ \int_{0}^{1}\int_{y}^{1} e^{x^{2}} \:\: dx\:dy $$
\end{tarea}



\begin{definition}  
    \label{def:area}
    Sea $D\subset\Rn{2}$,  se define el \'area de $D$,  y se denota $\text{A}(D)$,  a  
    \[
        \text{A}(D)=\iint_D1\:dA.
    \]
\end{definition}
